\section{Background and Motivation}
\label{sec:background}

\subsection{System setting}

We consider a rack-scale pool composed of compute nodes and memory appliances.
An appliance exports fixed-size 2~MB \emph{extents}; a compute node may map
4~KB pages within an extent. The data path is either CXL-style load/store or
one-sided RDMA. As in prior far-memory work~\cite{aguilera2018remote,
amaro2020fastswap}, the CPU on the memory appliance is not involved in ordinary
accesses. A replicated pool manager assigns extents, issues access tokens, and
balances capacity. It is not on the load/store path.

Our fault model allows crash-stop compute nodes, manager replicas, and memory
appliances; messages may be delayed, reordered, or duplicated. An appliance's
DRAM is lost on power failure, but its small nonvolatile generation table
survives restart. The network may partition and later heal. We assume a
majority of manager replicas remains available and use crash-tolerant consensus
for manager metadata~\cite{ongaro2014raft}. Byzantine behavior and malicious
firmware are outside our model. Applications are responsible for durability of
their data; \system protects ownership and isolation, not application contents.

\subsection{The reconfiguration gap}

Consider moving extent $e$ from compute node $A$ to $B$. Four actions are
necessary: stop $A$, copy live pages, authorize $B$, and publish $B$'s mapping.
No order makes crashes harmless. Revoking $A$ first can make the only live copy
unavailable. Publishing $B$ first permits concurrent writers. Updating only
the manager first leaves hardware enforcing old authority; updating hardware
first leaves durable metadata describing the old state.

This interval is the \emph{reconfiguration gap}. It appears in more operations
than migration: allocating an extent from a free list, reclaiming memory from a
failed node, changing tenant keys, and evacuating an appliance all cross it.
Traditional distributed file systems close related gaps with locks and logs
~\cite{thekkath1997frangipani}, while persistent-memory systems order data and
metadata writes within one machine~\cite{volos2011mnemosyne,xu2016nova}. A
disaggregated-memory mapping spans independent machines and a hardware access
control boundary, so neither technique alone is sufficient.

\subsection{Why scanning and synchronous logging fall short}

Our measurements of a straightforward rebuild illustrate the scale mismatch.
Reading a 16-byte owner record for every 2~MB extent in a 16~TB pool takes
11.8~s with batched RDMA; validating mappings at 32 compute nodes raises this
to 16.4~s. The time is paid even if the failed manager interrupted no operation.
Conversely, synchronously appending a replicated record for every 4~KB mapping
raises the median fault-service time from 7.3 to 31.6~$\mu$s and saturates the
leader near 2.1 million mappings/s. The result is a false choice between slow
recovery and an expensive common path.

\paragraph{Design goals.}
\system therefore targets four properties: (1) exclusive ownership despite
arbitrary crash points; (2) recovery proportional to incomplete operations;
(3) no manager or consensus traffic on ordinary memory accesses; and (4)
incremental deployment on appliances with a small durable table. It does not
make volatile remote memory persistent, mask application bugs, or provide
transactions across application objects.
